\documentclass[shownotes, xcolor = table]{beamer}

\usepackage{xifthen}
\usepackage{ifthen}
\usepackage{xstring}

\usepackage{verbatim}
\usepackage{listings}

\usepackage{algorithm}
% \usepackage{algorithmic}
\usepackage[noend]{algpseudocode}
\algrenewcommand{\algorithmiccomment}[1]{\hfill$\triangleright$ \textcolor{teal}{#1}}

\usetheme{CambridgeUS} % try Madrid, CambridgeUS
\usecolortheme{beaver} % try beaver, dolphin, seahorse
\usefonttheme{serif}
% \usefonttheme[onlymath]{serif} % try professionalfonts, serif

\usepackage{amsmath, amsfonts, amssymb, mathtools, pifont}
\newcommand{\cmark}{\ding{51}}%
\newcommand{\xmark}{\ding{55}}%
\def\checkmark{\tikz\fill[scale=0.5](0,.35) -- (.25,0) -- (1,.7) -- (.25,.15) -- cycle;} 

\usepackage{graphicx, subcaption}
\usepackage[export]{adjustbox}

% input: sequence of numbers, separated by commas
% output: sequence of numbers, separated by \quad
% See http://tex.stackexchange.com/a/60026/23098
% \usepackage{xstring}
\newcommand{\seq}[1]{%
    \noexpandarg    
    \StrSubstitute{#1}{,}{\;}[\ele]
    \ele}

% needed: \usepackage[export]{adjustbox}
\newcommand{\importikznocaption}[3]{% #1: width, #2: height, #3: tex
\begin{figure}[h!]
    \centering
    \begin{adjustbox}{max totalsize = {#1}{#2}, center}
	  \input{#3}
    \end{adjustbox}
  \end{figure}
}

\usepackage{wasysym}

\usepackage[normalem]{ulem} % strike through text
\newcommand{\soutthick}[1]{%
    \renewcommand{\ULthickness}{2.0pt}%
       \sout{#1}%
    \renewcommand{\ULthickness}{.4pt}% Resetting to ulem default
}

\setbeamersize{text margin left = 2em, text margin right = 1em}
\setbeamercolor{footnote mark}{fg = teal}
\setbeamertemplate{itemize items}[square]
% \setbeamertemplate{itemize items}[default]
\setbeamertemplate{enumerate items}[default]
\newcommand\coloreditem[1]{\item[\textcolor{#1}{\usebeamertemplate{itemize \beameritemnestingprefix item}}]}

\usepackage{tikz}
\usetikzlibrary{mindmap, arrows.meta, shapes, positioning, calc, backgrounds, fit}

\theoremstyle{plain}

\lstdefinestyle{code}{
  language = Java,
  basicstyle = \rmfamily\footnotesize,
  captionpos = t,
  mathescape = true,
  showstringspaces = false,
  keywordstyle = \color{blue}\bfseries,
  commentstyle = \color{teal},
  stringstyle = \color{green}\bfseries,
  numbers = left,
  numberstyle = \tiny,
  % upquote = true, 
  frame = lines,
  breaklines = true
}

% for tables
\usepackage{multirow}
\newcommand{\innercell}[2]{\begin{tabular}{@{}#1@{}}#2\end{tabular}}
\usepackage{hhline}
%%%%%%%%%%%%%% for appendix %%%%%%%%%%%%%%%%
% http://www-ljk.imag.fr/membres/Jerome.Lelong/latex/appendixnumberbeamer.sty
% Reference: http://tex.stackexchange.com/questions/2541/beamer-frame-numbering-in-appendix
\usepackage{appendixnumberbeamer}
% Add total frame count to slides, optional. From Stefan,
% http://www.latex-community.org/forum/viewtopic.php?f=4&t=2173
\expandafter\def\expandafter\insertshorttitle\expandafter{%
  \insertshorttitle\hfill\insertframenumber\,/\,\inserttotalframenumber}
%%%%%%%%%%%%%% for appendix %%%%%%%%%%%%%%%%

% for fig without caption: #1: width/size; #2: fig file
\newcommand{\fignocaption}[2]{
  \begin{figure}[htp]
    \centering
    \includegraphics[#1]{#2}
  \end{figure}
}

% for fig without caption: #1: width/size; #2: fig file; #3: fig caption
\newcommand{\fig}[3]{
  \begin{figure}[htp]
    \centering
      \includegraphics[#1]{#2}
      \caption{#3}
  \end{figure}
}

% for cite: #1: author; #2: conference #3: year
\newcommand{\citeinbeamer}[3]{{\tiny{\textcolor{blue}{[#1@#2'#3]}}}}

\usepackage[backend=bibtex]{biblatex}
\addbibresource{cs-application-report.bib}

\newcommand{\term}[1]{\textcolor{blue}{\scriptsize (#1)}}
\newcommand{\set}[1]{\{#1\}}
\newcommand{\question}[1]{\textcolor{red}{\centerline{#1}}}
\newcommand{\answer}[1]{\textcolor{blue}{\centerline {#1}}}
\newcommand{\alertred}[1]{\textcolor{red}{#1}}
\newcommand{\alertblue}[1]{\textcolor{blue}{#1}}
\newcommand{\todo}[1]{\textcolor{red}{\textbf{TODO:} #1}}
\newcommand{\mathbfblue}[1]{\textcolor{blue}{$\mathbf{#1}$}}

\newcommand{\reporttitle}{Bubble Sort}
%%%%%%%%%%%%%%%%%%%%%%%%%%%%%%%%%%%%%%%%%%%%%%%%%%%%%%%%%%%%%%%%%%%%%%%%%%%%%%%%	
\title[\reporttitle]{\reporttitle}
\subtitle{(What are Algorithms and How to Analyze Algorithms)}

\author[Hengfeng Wei]{Hengfeng Wei}
\titlegraphic{\includegraphics[height = 1.2cm]{figs/nju-logo-purple.png}~\includegraphics[height = 1.2cm]{figs/cs-logo.jpg}}
\institute[ICS@NJU]{Institute of Computer Software\\Nanjing University}
\date{\today}

\AtBeginSection[]{
  \begin{frame}[noframenumbering, plain]
    \frametitle{\reporttitle}
    \tableofcontents[currentsection, sectionstyle=show/shaded, subsectionstyle=hide/hide/hide]
  \end{frame}
}
%%%%%%%%%%%%%%%%%%%%
\begin{document}

\maketitle

\begin{frame}[noframenumbering, plain]
  \frametitle{\reporttitle}
  \tableofcontents[currentsection, sectionstyle=show, subsectionstyle=show/show/hide]
\end{frame}

\section{The Sorting Problem}

%%%%%%%%%%%%%%%%%%%%
\begin{frame}{Algorithms}
  % \begin{center}
  %   An algorithm is a sequence of operations \\
  %   that transform the input into the output.
  % \end{center}
  \begin{center}
	What is an algorithm? \qquad \pause What is computation?
  \end{center}

  \pause

  \fignocaption{width = 0.50\textwidth}{figs/algorithm-def.pdf}

  \pause
  \vspace{-0.60cm}
  \begin{center}
	\textcolor{red}{Correctness!}
  \end{center}

  \pause

  \begin{description}[Effectiveness:]
	\item[Definiteness:] precisely defined operations
	  \pause
	\item[Finiteness:] termination
	  \pause
	\item[Effectiveness:] a reasonable model; basic operations % RAM {\scriptsize (Random-Access Machine)} model
	  \pause
	  \begin{itemize}
		% \item unrealistic: \texttt{sort} operation
		% \item realistic: arithmetic, data movement, and control
		%   \pause
		\item for sorting: compare, swap
	  \end{itemize}
  \end{description}
\end{frame}
%%%%%%%%%%%%%%%%%%%%
\begin{frame}{Sorting}
  The sorting problem:
  \begin{description}
	\item[Input:] A sequence of $n$ integers $A$:\seq{$a_1, a_2, \cdots, a_n$}.
	\item[Output:] A permutation \seq{$a'_1, a'_2, \ldots, a'_n$} of $A$ \emph{s.t.} $a'_1 \le a'_2 \le \cdots \le a'_n$ {\small (non-decreasing order)}.
  \end{description}

  \[
	3\quad 1\quad 4\quad 2\quad \Longrightarrow 1\quad 2\quad 3\quad 4
  \]

  % sortable
  % A little more formalism: ordering relation ``$<$'' on $A$.

  % \vspace{0.20cm}
  % $\forall a, b, c \in A$,
  % \begin{description}[Transitivity:]
  %   \item[Trichotomy:] $a < b, a = b, a > b$
  %   \item[Transitivity:] $a < b \land b < c \implies a < c$
  % \end{description}
\end{frame}
%%%%%%%%%%%%%%%%%%%%
\begin{frame}{Inversions}
  \[
	A = a_1\quad a_2\quad \ldots\quad a_n.
  \]

  \begin{center}
	If $i < j$ and $a_{i} > a_{j}$, then $(a_i, a_j)$ is an \textcolor{red}{\bf inversion}.\\[8pt] \pause
	\textcolor{blue}{Adjacent} inversion: $(a_i, a_{i+1})$
  \end{center}

  \pause
  \vspace{-0.50cm}

  \begin{columns}
	\column{0.50\textwidth}
	  \fignocaption{width = 0.50\textwidth}{figs/inversions-example.pdf}
	\column{0.50\textwidth}
	{\small
	  \begin{center}
		\#inversions = 3\\
		\#adjacent inversions = 2
	  \end{center}
	}
  \end{columns}

  \pause
  \begin{columns}
	\column{0.50\textwidth}
	  \fignocaption{width = 0.50\textwidth}{figs/inversions-example-nonincreasing.pdf}
	\column{0.50\textwidth}
	{\small
	  \begin{center}
		\#inversions = 3 + 2 + 1 = 6\\
		\#adjacent inversions = 3
	  \end{center}
	}
  \end{columns}

  \pause
  \begin{columns}
	\column{0.50\textwidth}
	  \fignocaption{width = 0.50\textwidth}{figs/inversions-example-nondecreasing.pdf}
	\column{0.50\textwidth}
	{\small
	  \begin{center}
		\#inversions = 0\\
		\#adjacent inversions = 0
	  \end{center}
	}
  \end{columns}
\end{frame}
%%%%%%%%%%%%%%%%%%%%
\begin{frame}{Inversions}
  \begin{center}
	\fbox{\textcolor{blue}{Theorem:} $A$ is sorted $\iff$ $A$ has no adjacent inversions.}
  \end{center}

  \begin{align*}
	\onslide<2->{A \text{ is sorted } &\Longrightarrow A \text{ has no inversions} \\}
	  \onslide<3->{&\Longrightarrow A \text{ has no adjacent inversions}.}
  \end{align*}

  \vspace{-0.50cm}

  \begin{align*}
	\onslide<4->{A \text{ has no adjacent inversions } &\Longrightarrow \forall i \in [1,n-1]: a_{i} \le a_{i+1} \\}
	  \onslide<5->{&\Longrightarrow A \text{ is sorted}.}
  \end{align*}
\end{frame}
%%%%%%%%%%%%%%%%%%%%

\section{Bubble Sort}

%%%%%%%%%%%%%%%%%%%%
\begin{frame}{Bubble Sort}
  Basic idea: to eliminate all adjacent inversions
  \pause

  % basic idea of bubble sort with overlay

\begin{algorithm}[H]
  \begin{algorithmic}[1]
	\Procedure{BubbleSort}{$A: a_1\; a_2\; \cdots\; a_n$}
	  \Repeat
		\State pick any $i$		
		\uncover<3->{\Comment{\textcolor{red}{\bf Definiteness!}}}
		  \If{$a_{i} > a_{i+1}$}	\Comment{CAS}
			\State \Call{Swap}{$a_{i}, a_{i+1}$}
		  \EndIf
	  \Until{no adjacent inversions}	
	  \uncover<3->{\Comment{\uncover<4->{\textcolor{red}{\bf Finiteness!}}
		\textcolor{red}{\bf Definiteness!}}}
    \EndProcedure
  \end{algorithmic}
\end{algorithm}

\end{frame}
%%%%%%%%%%%%%%%%%%%%
\begin{frame}{Bubble Sort: Definiteness}
  % bubble sort alg

\begin{algorithm}[H]
  \begin{algorithmic}[1]
	\Procedure{BubbleSort}{$A: a_1\; a_2\; \cdots\; a_n$}
	  \Repeat
		\State \fbox{\uncover<2->{swapped $\gets$ false}}
		\For{$i \gets 1 : n - 1$}	\Comment{Pick $i$}
		  \If{$a_{i} > a_{i+1}$}	% \Comment{CAS}
			\State \Call{Swap}{$a_{i}, a_{i+1}$}
			\State \fbox{\uncover<3->{swapped $\gets$ true}}
		  \EndIf
		\EndFor
		\Until{\fbox{\uncover<4->{swapped = false}}}
    \EndProcedure
  \end{algorithmic}
\end{algorithm}

\end{frame}
%%%%%%%%%%%%%%%%%%%%
\begin{frame}{Bubble Sort: Example}
  % \begin{center}
  %   \scalebox{0.50}{
  %     \begin{minipage}{0.70\linewidth}
  %   	% bubble sort

\begin{algorithm}[H]
  \begin{algorithmic}[1]
	\Procedure{BubbleSort}{$A: a_1\; a_2\; \cdots\; a_n$}
	  \Repeat
		\State swapped $\gets$ false
		\For{$i \gets 1 : n - 1$}	\Comment{Pick $i$}
		  \If{$a_{i} > a_{i+1}$}	% \Comment{CAS}
			\State \Call{Swap}{$a_{i}, a_{i+1}$}
			\State swapped $\gets$ true
		  \EndIf
		\EndFor
		\Until{swapped = false}		\Comment{No swaps}
    \EndProcedure
  \end{algorithmic}
\end{algorithm}

  %     \end{minipage}
  %   }
  % \end{center}

  \begin{center}
	\begin{tikzpicture}[rotate = 90, 
	swap/.style = {dashed, thick},
	inv/.style = {blue, rectangle, draw, inner sep = 0pt, outer sep = 0pt},
	ele/.style = {font = \Large},
	pass/.style = {brown, thick, dashed},
	bubble/.style = {violet, bend right = #1, decorate, decoration = {shape backgrounds, shape = circle, shape size = 5pt, shape sep = 12pt}}]
	\uncover<1->{
  \foreach \e [count = \ei] in {3,1,4,2} {
	  \node (0\e) [ele] at (\ei, 1) {$\e$};
  }
  \draw [pass] (0.5, 0.5) to (4.5,0.5);
}

\uncover<2->{
  \foreach \e [count = \ei] in {3,1,4,2} {
	  \node (1\e) [ele] at (\ei, 0) {\ifthenelse{\e = 3 \OR \e = 1}{\textcolor{red}{$\e$}}{$\e$}};
  }
  \begin{pgfonlayer}{background}
    \node<2-> () [inv, fit = (13) (11)] {};
  \end{pgfonlayer}
}

\uncover<3->{
  \foreach \e [count = \ei] in {1,3,4,2} {
    \node (2\e) [ele] at (\ei, -1) {$\e$};
  }
  \begin{pgfonlayer}{background}
    \node<3-> () [inv, fit = (23) (24)] {};
  \end{pgfonlayer}
}

\uncover<4->{
  \foreach \e [count = \ei] in {1,3,4,2} {
	  \node (3\e) [ele] at (\ei, -2) {\ifthenelse{\e = 4 \OR \e = 2}{\textcolor{red}{$\e$}}{$\e$}};
  }
  \begin{pgfonlayer}{background}
    \node<4-> () [inv, fit = (34) (32)] {};
  \end{pgfonlayer}
  % the second pass
  \draw [pass] (0.5, -2.5) to (4.5,-2.5);
}

\uncover<5->{
  \foreach \e [count = \ei] in {1,3,2,4} {
	  \node (4\e) [ele] at (\ei, -3) {$\e$};
  }
  \begin{pgfonlayer}{background}
    \node<5-> () [inv, fit = (41) (43)] {};
  \end{pgfonlayer}
}

\uncover<6->{
  \foreach \e [count = \ei] in {1,3,2,4} {
	  \node (5\e) [ele] at (\ei, -4) {\ifthenelse{\e = 3 \OR \e = 2}{\textcolor{red}{$\e$}}{$\e$}};
  }
  \begin{pgfonlayer}{background}
    \node<6-> () [inv, fit = (53) (52)] {};
  \end{pgfonlayer}
}

\uncover<7->{
  \foreach \e [count = \ei] in {1,2,3,4} {
	  \node (6\e) [ele] at (\ei, -5) {$\e$};
  }
  \begin{pgfonlayer}{background}
    \node<7-> () [inv, fit = (63) (64)] {};
  \end{pgfonlayer}

  \draw [pass] (0.5, -5.5) to (4.5,-5.5);
}

\uncover<8->{
  \foreach \e [count = \ei] in {1,2,3,4} {
	  \node (7\e) [ele] at (\ei, -6) {$\e$};
  }
  \begin{pgfonlayer}{background}
    \node<8-> () [inv, fit = (71) (72)] {};
  \end{pgfonlayer}
}

\uncover<9->{
  \foreach \e [count = \ei] in {1,2,3,4} {
	  \node (8\e) [ele] at (\ei, -7) {$\e$};
  }
  \begin{pgfonlayer}{background}
    \node<9-> () [inv, fit = (82) (83)] {};
  \end{pgfonlayer}
}

\uncover<10->{
  \foreach \e [count = \ei] in {1,2,3,4} {
	  \node (9\e) [ele] at (\ei, -8) {$\e$};
  }
  \begin{pgfonlayer}{background}
    \node<10-> () [inv, fit = (93) (94)] {};
  \end{pgfonlayer}
}

\uncover<11->{
  \draw [bubble = 15] (04) to (44);
}

\uncover<12->{
  \draw [bubble = 15] (33) to (73);
}
\end{tikzpicture}

  \end{center}

  % \[
  %   5\quad 3\quad 1\quad 6\quad 7\quad 2\quad 4\quad 8
  % \]
\end{frame}
%%%%%%%%%%%%%%%%%%%%
\begin{frame}{Bubble Sort: Finiteness}
  \begin{center}
	\scalebox{0.75}{
	  \begin{minipage}{0.70\linewidth}
		% bubble sort

\begin{algorithm}[H]
  \begin{algorithmic}[1]
	\Procedure{BubbleSort}{$A: a_1\; a_2\; \cdots\; a_n$}
	  \Repeat
		\State swapped $\gets$ false
		\For{$i \gets 1 : n - 1$}	\Comment{Pick $i$}
		  \If{$a_{i} > a_{i+1}$}	% \Comment{CAS}
			\State \Call{Swap}{$a_{i}, a_{i+1}$}
			\State swapped $\gets$ true
		  \EndIf
		\EndFor
		\Until{swapped = false}		\Comment{No swaps}
    \EndProcedure
  \end{algorithmic}
\end{algorithm}

	  \end{minipage}
	}
  \end{center}

  The inner ``{\bf for}'' loops:
  \begin{enumerate}[1)]
	\item \textcolor{blue}{$\exists$ loop : no swaps} $\implies$ swapped = false $\implies$ terminates
	  \pause
	\item \textcolor{red}{$\forall$ loop : has swaps} \pause \qquad \textcolor{red}{\bf \textsc{Impossible}!}
  \end{enumerate}
\end{frame}
%%%%%%%%%%%%%%%%%%%%
\begin{frame}{Bubble Sort: Finiteness}
  \begin{center}
	\fbox{Total \#inversions is finite.}
  \end{center}

  \pause
  Effects of \textsc{Swap}{($a_{i}, a_{i+1}$)} on \textcolor{red}{\#inversions}:
  \fignocaption{width = 0.45\textwidth}{figs/swap-inversions.pdf}

  \pause
  \begin{center}
	\textcolor{blue}{$-1: (a_{i},a_{i+1})$} \pause \qquad\qquad \textcolor{red}{$+0$}
  \end{center}

  \pause
  \begin{center}
	\fbox{$\textsc{Swap}(a_{i}, a_{i+1}) \implies -1 \text{ inversion}$}
  \end{center}
\end{frame}
%%%%%%%%%%%%%%%%%%%%
\begin{frame}{Bubble Sort: Correctness}
  \begin{align*}
	\onslide<1->{\text{\textcolor{blue}{\bf Finiteness }}}
	  \onslide<2->{&\implies \exists \text{ loop : no swaps} \\[3pt]}
	  \onslide<3->{&\implies A \text{ has no adjacent inversions any more} \\[3pt]}
	  \onslide<4->{&\implies A \text{ is already sorted.}}
  \end{align*}
\end{frame}
%%%%%%%%%%%%%%%%%%%%
\begin{frame}{Optimizing Bubble Sort (I)\footnote{See Appendix for ``Optimizing Bubble Sort (II)''.}}
  Idea: After each ``{\bf for}'' loop, one more element is settled.

  % mfp: maximal in its final position

\begin{algorithm}[H]
  \begin{algorithmic}[1]
	\Procedure{BubbleSort}{$A: a_1\; a_2\; \cdots\; a_n$}
	  \State $n \gets \text{len}(A)$
	  \Repeat
		\State swapped $\gets$ false
		\For{$i \gets 1 : n - 1$}	% \Comment{Pick $i$}
		  \If{$a_{i} > a_{i+1}$}	% \Comment{CAS}
			\State \Call{Swap}{$a_{i}, a_{i+1}$}
			\State swapped $\gets$ true
		  \EndIf
		\EndFor
		\State $n \gets n - 1$	\Comment{One maximal bubbles up}
		\Until{swapped = false}
    \EndProcedure
  \end{algorithmic}
\end{algorithm}

\end{frame}
%%%%%%%%%%%%%%%%%%%%
\section{Analysis of Bubble Sort}

%%%%%%%%%%%%%%%%%%%%
\begin{frame}{Time Complexity of Bubble Sort}
  \begin{itemize}
	\setlength{\itemsep}{10pt}
	\item Finiteness is NOT enough $\implies$ Quantifying finiteness
	  \pause
	\item Time on real computers varies $\implies$ \#Ops on our model:
	  \vspace{5pt}
	  \pause
	  \begin{description}
		\setlength{\itemsep}{5pt}
		% \item[$|P|:$] \#Passes \hfill (the ``{\bf for}'' loops)
		\item[$|C|:$] \#Comparisons \hfill ({\bf if} $a_{i} > a_{i+1}$)
		\item[$|S|:$] \#Swaps  \hfill ($\textsc{Swap}(a_{i}, a_{i+1})$)
	  \end{description}
	  % \pause
	  % \[
	  %   |C| \ge |S|
	  % \]
	  % \vspace{-0.50cm}
	  \pause
	\item Different inputs $\implies$ $|C|$ and $|S|$ vary:
	  \vspace{5pt}
	  \pause
	  \begin{itemize}
		\item Best-case, worst-case, and average-case analysis
	  \end{itemize}
\end{itemize}
\end{frame}
%%%%%%%%%%%%%%%%%%%%
\begin{frame}{Best-case and Worst-case Analysis}
  \uncover<2->{
  \[
	\text{\textcolor{blue}{Best-case:} } 1\quad 2\; \cdots\; n\quad
  \]}
  
  \begin{displaymath}
	\begin{array}{lll}
	  & \hfill \textcolor{blue}{\frac{\text{Best-case:}}{\uncover<2->{\text{non-decreasingly sorted}}}} \hfill 
	  & \hfill \textcolor{red}{\frac{\text{Worst-case:}}{\uncover<4->{\text{non-increasingly sorted}}}} \\[10pt]
	  % |P| & = (\uncover<3->{\min: 1,} & \uncover<5->{\max: n}); \\[6pt]
	  |C| & = (\uncover<3->{\min: n-1,} & \uncover<5->{\max: \frac{n^2 - n}{2}}); \\[6pt]
	  |S| & = (\uncover<3->{\min: 0,} & \uncover<5->{\max: \frac{n^2 - n}{2}}).
	\end{array}
  \end{displaymath}

  \uncover<4->{
	\[
	  \text{\textcolor{red}{Worst-case: }} n\quad n-1\; \cdots\; 1
	\]}
  % \uncover<5->{
  %   \[
  %     \#\text{inversions} = (n-1) + (n-2) + \cdots + 1 = \frac{n^2 - n}{2}.
  %   \]
  % }
\end{frame}
%%%%%%%%%%%%%%%%%%%%
\begin{frame}{$|S|:$ \#Swaps (Average-case Analysis)\footnote{An exercise: what is $|C|$ (\#Comparisions) in average?}}
  Assumptions on inputs:
  \begin{enumerate}
	\item The input is a random permutation (``average input'')
	  \pause
	\item All numbers are different (for simplicity)
  \end{enumerate}

  \pause

  \begin{center}
	\fbox{\textcolor{blue}{Lemma:} $\textsc{Swap}(a_{i}, a_{i+1}) \implies -1 \text{ inversion}$}
  \end{center}

  \[
	|S| = \textcolor{red}{\mathbb{E}(\#\text{inversions})}
  \]
\end{frame}
%%%%%%%%%%%%%%%%%%%%
\begin{frame}{$|S|:$ \#Swaps (Average Analysis)}
  \begin{align*}
	% \onslide<1->{I_{ij} &: \text{indicator of inversion } (a_i, a_j) \\[3pt]}
	\onslide<1->{I_{ij} &= \left\{\begin{array}{lr}
		1 & (a_i, a_j) \text{ is an inversion}\\
		0 & \text{ o.w.} 
	  \end{array}\right.\\[3pt]}
	\onslide<1->{X &= \sum_{1 \le i < n} \sum_{i < j \le n} I_{ij}\qquad\quad (\text{\#inversions}) \\[3pt]}
	\onslide<2->{\mathbb{E}(X) &= \mathbb{E}(\sum_{i} \sum_{j>i} I_{ij})}\onslide<3->{= \sum_{i} \sum_{j>i} \mathbb{E}(I_{ij})\quad (\text{linearity of expectation})\\[3pt]}
	\onslide<4->{\mathbb{E}(I_{ij}) &= \mathbb{P}\set{I_{ij} = 1}} \onslide<5->{{}={} \frac{1}{2}\qquad (a_i \neq a_j; \text{half: } a_i < a_j, \text{half: } a_i > a_j) \\[3pt]}
	\onslide<6->{\mathbb{E}(X) &= \sum_{i} \sum_{i<j} \frac{1}{2} = \binom{n}{2} \cdot \frac{1}{2} = \frac{n(n-1)}{4} \textcolor{red}{{}={} O(n^2)}}
  \end{align*}
\end{frame}
%%%%%%%%%%%%%%%%%%%%
\begin{frame}{Faster Algorithms}
  \begin{quote}
	It took a good deal of work to analyze the bubble sort;
	and although [\dots], 
	the results are disappointing 
	since they tell us that \textcolor{red}{the bubble sort isn't really very good at all}.\\
	\hfill --- Donald E. Knuth
  \end{quote}

  \pause

  \begin{center}
	faster: $O(n^2) \to O(n \lg n)$?\\[5pt] \pause
	\dots and faster: $O(n \lg n) \to O(n)$?
  \end{center}

  \pause
  \fignocaption{width = 0.18\textwidth}{figs/stay-tuned.jpg}
\end{frame}
%%%%%%%%%%%%%%%%%%%%
\begin{frame}[noframenumbering]
  \fignocaption{width = 0.20\textwidth}{figs/qa.png}
  \vspace{-0.8cm}
  \begin{center}
    \textcolor{blue}{\bf \large hengxin0912@gmail.com}
  \end{center}
  \vspace{-0.5cm}
  \fignocaption{width = 0.50\textwidth}{figs/thankyou.jpg}
\end{frame}
%%%%%%%%%%%%%%%%%%%%


\end{document}